\documentclass{article}
\usepackage[T1]{fontenc}
\usepackage{lmodern}
\usepackage{graphicx}
\usepackage{amsmath,amssymb}
\usepackage[a4paper,margin=2cm]{geometry}
\usepackage[absolute,overlay]{textpos}
\setlength{\TPHorizModule}{1mm}
\setlength{\TPVertModule}{1mm}
\usepackage[indonesian]{babel}
\usepackage{hyperref}
\setlength{\emergencystretch}{2em}
% unit: Halaman 24

\begin{document}

% === Halaman 24 (llm) ===

\section*{Man-in-the-middle Attack}

\begin{itemize}
    \item Karena pengirim dan penerima harus berbagi kunci publik, maka distribusi kunci publik dapat mengalami serangan \textit{man-in-the-middle attack}.
    \item Misalkan Alice dan Bob mengirim kunci publiknya masing-masing melalui saluran komunikasi. Orang-di-tengah, misalkan Carol, mengintersepsi komunikasi antara Bob dan Alice lalu ia berpura-pura sebagai salah satu pihak (Alice atau Bob).
    \item Carol (yang menyamar sebagai Alice) mengirim kunci publiknya kepada Bob (Bob percaya itu adalah kunci publik Alice), dan Carol (yang menyamar sebagai Bob) mengirim kunci publiknya kepada Alice (Alice percaya itu adalah kunci publik Bob).
    \item Selanjutnya, Carol mendekripsi pesan dari Bob dengan kunci privatnya, menyimpan salinannya, lalu mengenkripsi pesan tersebut dengan kunci publik Alice, dan mengirim cipherteks tersebut kepada Alice. Alice dan Bob tidak dapat mendeteksi keberadaan Carol.
\end{itemize}

\end{document}
